% !TeX document-id = {df9ec23f-6055-43ee-aadd-d039429067b8}
% !TeX TXS-program:compile = txs:///latexmk/{-pdf}
% !TeX root = main.tex
% Kompilacja: latexmk main.tex (silnik LuaLaTeX wybiera latexmkrc)
% Pełna instrukcja konfiguracji TeXstudio znajduje się w pliku README.md.

% === KONFIGURACJA SZABLONU — NIE ZMIENIAJ W TYM PLIKU ===
% Dane prezentacji, język, typ pracy, logotypy i ostatni slajd ustawia się
% wyłącznie w config/metadata.tex.
% JEDYNY PLIK, KTORY TRZEBA EDYTOWAC PRZED PRZYGOTOWANIEM PREZENTACJI.
% Jezyk: polish lub english. Typ pracy: engineering, master lub doctoral.
\newcommand{\presentationlanguage}{polish}
\newcommand{\thesistype}{engineering}

\newcommand{\universitypl}{Politechnika Opolska}
\newcommand{\universityen}{Opole University of Technology}
\newcommand{\facultypl}{Wydział Elektrotechniki, Automatyki i Informatyki}
\newcommand{\facultyen}{Faculty of Electrical Engineering, Automatic Control and Informatics}

\newcommand{\titlepl}{Tytuł pracy dyplomowej}
\newcommand{\titleen}{Diploma thesis title}
\newcommand{\authorname}{Imię i nazwisko}
\newcommand{\albumid}{000000}
\newcommand{\programmepl}{Nazwa kierunku}
\newcommand{\programmeen}{Study programme}
\newcommand{\supervisorname}{dr hab. inż. Imię Nazwisko, prof. uczelni}
\newcommand{\auxsupervisorname}{}
\newcommand{\city}{Opole}
\newcommand{\presentationyear}{2026}

% Logo KreativEU mozna wylaczyc, ustawiajac false.
\newcommand{\showkreativeulogo}{true}

% Ostatni slajd: title (powtorzenie strony tytulowej), thanks albo none.
\newcommand{\lastslidetype}{title}

% Pliki identyfikacji wizualnej.
\newcommand{\universitylogopl}{assets/branding/logo4-Politechnika_Opolska.pdf}
\newcommand{\universitylogoen}{assets/branding/logo-en-cropped.pdf}
\newcommand{\facultylogo}{assets/branding/Logo-WEAIL.png}
\newcommand{\kreativeulogo}{assets/branding/logos-KreativEU-vector-color.png}

\documentclass[aspectratio=169,11pt]{beamer}
\usepackage{config/presentation-style}

\title{\presentationtitle}
\author{\authorname}
\date{\presentationyear}

\begin{document}

% === SLAJDY POCZĄTKOWE — NIE ZMIENIAJ KOLEJNOŚCI POLECEŃ ===
% Tytuł, autor i pozostałe dane są pobierane z config/metadata.tex.
\begin{frame}
  \titlepage
\end{frame}

\begin{frame}{\agendatitle}
  \tableofcontents[hideallsubsections]
\end{frame}

% === SEKCJE PREZENTACJI — TUTAJ USTALASZ ICH KOLEJNOŚĆ ===
% Każdą część umieść w osobnym pliku w katalogu sections/ i dołącz przez
% \input{sections/nazwa-pliku}, bez rozszerzenia .tex.
% Nową strukturę można utworzyć poleceniem:
% python tools/create_section.py 04-nazwa --title-pl "Tytuł" --title-en "Title"
% Polecenie \section ustala nazwę w spisie treści; powinno bezpośrednio
% poprzedzać odpowiadające mu polecenie \input.
\section{\problemsectiontitle}
\begin{frame}{Problem, który rozwiązuje praca}
  \begin{columns}[T,onlytextwidth]
    \begin{column}{0.54\textwidth}
      \begin{itemize}
        \item Krótko opisz kontekst i praktyczny problem.
        \item Wskaż lukę w istniejących rozwiązaniach.
        \item Wyjaśnij, dlaczego problem jest istotny.
      \end{itemize}
      \vspace{3mm}
      \begin{block}{Główna teza}
        Jedno zdanie, które słuchacze powinni zapamiętać z tego slajdu.
      \end{block}
    \end{column}
    \begin{column}{0.42\textwidth}
      \imageplaceholder[0.38\textheight]{Miejsce na rysunek problemu, fotografię stanowiska albo schemat systemu}
    \end{column}
  \end{columns}
\end{frame}

\begin{frame}{Cel i zakres pracy}
  \begin{columns}[T,onlytextwidth]
    \begin{column}{0.48\textwidth}
      \textbf{Cel główny}
      \begin{itemize}
        \item Zdefiniuj mierzalny rezultat pracy.
      \end{itemize}
      \vspace{3mm}
      \textbf{Zakres}
      \begin{itemize}
        \item analiza wymagań,
        \item projekt i implementacja,
        \item badania oraz ocena wyników.
      \end{itemize}
    \end{column}
    \begin{column}{0.48\textwidth}
      \textbf{Kryteria sukcesu}
      \begin{enumerate}
        \item Kryterium jakościowe lub ilościowe.
        \item Warunek brzegowy albo ograniczenie.
        \item Sposób weryfikacji rezultatu.
      \end{enumerate}
    \end{column}
  \end{columns}
\end{frame}


\section{\methodssectiontitle}
\begin{frame}{Przyjęta metodyka prowadzi od wymagań do weryfikacji}
  \begin{enumerate}
    \item \highlight{Analiza} -- wymagania, dane wejściowe i ograniczenia.
    \item \highlight{Projekt} -- architektura rozwiązania i dobór narzędzi.
    \item \highlight{Realizacja} -- implementacja oraz integracja elementów.
    \item \highlight{Walidacja} -- scenariusze badań i kryteria oceny.
  \end{enumerate}
  \vspace{5mm}
  \imageplaceholder[0.22\textheight]{Miejsce na schemat procesu lub architekturę rozwiązania}
\end{frame}

\begin{frame}{Najważniejszy element rozwiązania}
  \begin{columns}[T,onlytextwidth]
    \begin{column}{0.57\textwidth}
      \imageplaceholder[0.43\textheight]{Miejsce na diagram, model, zrzut interfejsu albo fotografię prototypu}
    \end{column}
    \begin{column}{0.39\textwidth}
      \begin{itemize}
        \item Nazwij element i jego rolę.
        \item Wyjaśnij przepływ informacji.
        \item Wskaż kluczową decyzję projektową.
        \item Podaj najważniejsze ograniczenie.
      \end{itemize}
    \end{column}
  \end{columns}
\end{frame}

\begin{frame}[fragile]{Przykład kodu źródłowego}
\begin{lstlisting}[style=weail,language=Python]
def moving_average(samples, window):
    """Oblicz średnią kroczącą z próbek."""
    # Obsługa polskich znaków: ąćęłńóśźż
    return [
        sum(samples[i:i + window]) / window
        for i in range(len(samples) - window + 1)
    ]
\end{lstlisting}
  \vspace{2mm}
  \small Pokazuj wyłącznie fragment potrzebny do uzasadnienia decyzji technicznej.
\end{frame}


\section{\resultssectiontitle}
\begin{frame}{Wyniki potwierdzają osiągnięcie celu}
  \begin{center}
    \metric{98\%}{skuteczność rozwiązania}\hfill
    \metric{35 ms}{czas odpowiedzi}\hfill
    \metric{12\%}{poprawa względem wariantu bazowego}
  \end{center}
  \vspace{7mm}
  \imageplaceholder[0.25\textheight]{Miejsce na wykres z jednostkami, legendą i podpisanymi osiami}
\end{frame}

\begin{frame}{Porównanie wariantów}
  \centering
  \small
  \begin{tabularx}{\textwidth}{Xrrr}
    \toprule
    \textbf{Wariant} & \textbf{Jakość} & \textbf{Czas} & \textbf{Koszt} \\
    \midrule
    Bazowy & 0,82 & 52 ms & niski \\
    Proponowany & \textbf{0,94} & \textbf{35 ms} & średni \\
    Referencyjny & 0,91 & 41 ms & wysoki \\
    \bottomrule
  \end{tabularx}
  \vspace{6mm}
  \begin{block}{Interpretacja}
    Nie poprzestawaj na tabeli: dodaj jedno zdanie wyjaśniające, co wynik zmienia w ocenie rozwiązania.
  \end{block}
\end{frame}

\begin{frame}{Wnioski i dalsze prace}
  \begin{columns}[T,onlytextwidth]
    \begin{column}{0.48\textwidth}
      \textbf{Najważniejsze wnioski}
      \begin{itemize}
        \item Cel pracy został osiągnięty w zdefiniowanym zakresie.
        \item Największy wpływ na wynik miał wskazany element.
        \item Rozwiązanie spełnia przyjęte kryteria.
      \end{itemize}
    \end{column}
    \begin{column}{0.48\textwidth}
      \textbf{Dalszy rozwój}
      \begin{itemize}
        \item rozszerzenie zbioru danych lub testów,
        \item optymalizacja wydajności,
        \item wdrożenie w środowisku docelowym.
      \end{itemize}
    \end{column}
  \end{columns}
\end{frame}


% === BIBLIOGRAFIA — NIE ZMIENIAJ SPOSOBU JEJ GENEROWANIA ===
% Źródła dodawaj do bibliografia.bib i cytuj w slajdach przez \cite{klucz}.
% Usuń \nocite{*}, jeżeli bibliografia ma zawierać wyłącznie cytowane pozycje.
\begin{frame}[allowframebreaks]{\bibliographytitle}
  \nocite{*}
  \printbibliography[heading=none]
\end{frame}

% === OSTATNI SLAJD — NIE EDYTUJ PONIŻSZEJ LOGIKI ===
% Wybierz wariant title, thanks albo none przez \lastslidetype w
% config/metadata.tex.
\ifdefstring{\lastslidetype}{title}{%
  \begin{frame}
    \titlepage
  \end{frame}
}{%
  \ifdefstring{\lastslidetype}{thanks}{%
    \begin{frame}
      \vspace*{20mm}
      \begin{center}
        {\color{weailblue}\fontsize{28}{32}\selectfont\bfseries \thankstitle}\par
        \vspace{5mm}
        {\large \questionslabel}\par
        \vspace{12mm}
        {\small\color{weailmuted}\authorname\quad | \quad\presentationtitle}
      \end{center}
    \end{frame}
  }{%
    \ifdefstring{\lastslidetype}{none}{}{%
      \PackageError{main}{Nieznana wartosc lastslidetype: \lastslidetype}{Uzyj title, thanks albo none.}%
    }%
  }%
}

\end{document}
